\section{Examples}\label{sec:problem-examples}

The preceding sections give detailed descriptions of every data block that can appear
in a problem file, but they do not make it obvious how the blocks are assembled to
calculate complete trajectories. This section therefore presents four example
problems and illustrates several TAOS capabilities.

The first example calculates a set of simple ballistic-reentry trajectories and shows
how surveys and summary variables are used for tradeoff studies. Although the
problem is simple, it shows the overall structure of a problem file. More complicated
trajectories use the same basic organization.

The second example calculates a rail-launched, two-stage ballistic rocket and uses
optimization to satisfy trajectory constraints. It illustrates multiple trajectories,
multiple segments, and a relatively simple optimization loop.

The third example computes a long-range intercept of a ballistic missile with
optimization. It is a more complex problem in which optimization shapes the
interceptor trajectory to meet the target.

The last example also computes an intercept, but it uses a TAOS intercept-guidance
algorithm instead of optimization.

The processing times quoted in these examples were measured on a Silicon Graphics
Indigo-2 workstation with a 200-MHz R4400 processor and are included only as
historical context for the 1995 implementation.
